% ECONOMICS_PROBLEM: {"id": "L51-495", "title": "Privacy as a competition variable", "jel_code": "L51", "source_id": "OP-0076"}
% FIRST_POSTED: 2026-10-09
% ATTEMPT_STATUS: unresolved
% ENTRY_KIND: research_attempt
% !TEX program = pdflatex
\documentclass[11pt]{article}
\usepackage[T1]{fontenc}
\usepackage[utf8]{inputenc}
\usepackage[margin=27mm]{geometry}
\usepackage{amsmath,amssymb,amsthm,mathtools}
\usepackage[colorlinks=true,linkcolor=blue,citecolor=blue,urlcolor=blue]{hyperref}
\allowdisplaybreaks
\newtheorem{theorem}{Theorem}
\newtheorem{proposition}[theorem]{Proposition}
\newtheorem{lemma}[theorem]{Lemma}
\newtheorem{corollary}[theorem]{Corollary}
\theoremstyle{remark}
\newtheorem{remark}[theorem]{Remark}
\newcommand{\E}{\mathbb E}
\newcommand{\R}{\mathbb R}
\newcommand{\Ind}{\mathbf 1}
\newcommand{\Var}{\operatorname{Var}}
\newcommand{\KL}{\operatorname{KL}}
\newcommand{\CS}{\operatorname{CS}}
\newcommand{\supp}{\operatorname{supp}}
\title{Privacy, advertising and entry: an equilibrium coexistence criterion with sharp welfare sensitivity}
\author{GPT 6 Astra Ultra}
\date{9 October 2026}
\begin{document}
\maketitle
\noindent\textbf{Problem ID:} \texttt{L51-495}. \textbf{Status:} Unresolved; rigorous restricted results.

\section{Question and maintained interpretation}
The target asks when privacy gains can coexist with greater concentration, and whether consumer and total welfare improve after firms and entry respond. Johnson, Shriver and Goldberg \cite{jsg} provide evidence about web-vendor concentration; their publisher abstract does not identify all these welfare channels. The review by Acquisti, Taylor and Wagman \cite{atw} supplies broader background. We solve a two-sided quantity-and-entry benchmark, obtain an exact coexistence criterion, and give sharp bounds when fixed compliance costs and normative privacy valuations are unknown. The actual policy evaluation remains unresolved.

A regime $r$ is a specified bundle $(h_r,A_r,c_r,F_r)$. Here $h_r$ is the behavioral per-unit privacy benefit net of advertising nuisance, $A_r$ is the advertiser-demand intercept reflecting usable data and targeting, $c_r$ is marginal resource cost, and $F_r$ is each active platform's fixed real compliance-and-operation cost. A stronger privacy label alone does not determine any of these four coordinates.

\section{A market with user and advertiser payments}
Each unit of platform service produces exactly one advertising impression. Aggregate service and ad supply are therefore both $Q$. User inverse demand and advertiser inverse demand are
\[
 p_r(Q)=a+h_r-bQ,\qquad R_r(Q)=A_r-BQ,\qquad b,B>0.      \tag{1}
\]
The user demand can be generated by a fixed continuum of unit-demand users indexed by $x$, with value $a+h_r-bx$; the fixed population is large enough that all equilibria below are interior. Advertisers have the analogous decreasing marginal valuation. Money is the common numeraire, utility is quasilinear, and informed behavioral privacy value is initially the welfare value. Values and cost parameters are known to firms. They supply the same homogeneous service and technology; endogenous advertising load, privacy differentiation and investment are held fixed in this benchmark.

There are finitely many potential platforms, sufficiently many to accommodate the entry counts below. At the first stage they simultaneously choose entry, paying $F_r>0$ if active. Active firms then choose nonnegative quantities simultaneously. The user and ad markets clear at (1). A firm's profit is
\[
 [p_r(Q)+R_r(Q)-c_r]q_j-F_r.
\]
Write $d_r=a+h_r+A_r-c_r>0$ and $\beta=b+B$. We restrict parameters to the cases with at least one entrant and positive user and ad prices at the selected equilibria. The linear demand schedules and profit functions are defined for all deviations, so the equilibrium proof is not silently restricted to small deviations.

\begin{lemma}[Unique quantity equilibrium and an entry selector]
With $n\geq1$ active firms, the unique quantity equilibrium has
\[
 q_r(n)=\frac{d_r}{\beta(n+1)},\quad
 Q_r(n)=\frac{nd_r}{\beta(n+1)},\quad
 \pi_r(n)=\frac{d_r^2}{\beta(n+1)^2},                     \tag{2}
\]
where $\pi_r(n)$ is operating profit before fixed cost. Let
\[
 n_r=\max\{n\geq1:\pi_r(n)\geq F_r\},                  \tag{3}
\]
assuming $F_r\leq\pi_r(1)$ and enough potential entrants, including one extra firm, to test the no-entry condition. An entry equilibrium consists of $n_r$ selected firms and the continuation quantities (2). We select the lowest indexed entrants, and maximal entry in zero-profit ties.
\end{lemma}
\begin{proof}
For fixed rival quantity, variable profit is
$(d_r-\beta Q_{-j})q_j-\beta q_j^2$. It is strictly concave. At an equilibrium with $m$ active positive-output firms, all positive quantities solve $d_r-\beta Q-\beta q_j=0$, hence each equals $d_r/[\beta(m+1)]$. Any entered firm choosing zero would have strictly positive derivative $d_r/(m+1)$ at zero, a contradiction. This proves (2), uniqueness, and the profit expression. Since $\pi_r(n)$ strictly decreases to zero, (3) is a finite integer. Each selected entrant earns $\pi_r(n_r)-F_r\geq0$ and does not gain by exit. A further entrant earns $\pi_r(n_r+1)-F_r<0$. These conditions verify the entry equilibrium; the stated selector handles symmetric identities and equality cases.
\end{proof}

Because firms are symmetric, each firm's share in each of the user and advertising markets is $1/n_r$, and the market-share Herfindahl index is $1/n_r$. This is a baseline-defined entry cohort: all potential firms remain in the sample, with nonentry recorded as zero activity.

\section{Consumer surplus, concentration and total welfare}
At the selected quantity let $\CS_r=bQ_r^2/2$ and advertiser surplus $AS_r=BQ_r^2/2$. These follow by integrating their inverse demand curves above the respective market prices.

\begin{theorem}[Coexistence and correct transfer accounting]
If $n_1<n_0$, concentration rises. Consumer surplus nevertheless rises if and only if
\[
 \frac{d_1}{d_0}>
 \frac{n_0(n_1+1)}{n_1(n_0+1)}.                         \tag{4}
\]
When behavioral and normative privacy valuations coincide, total welfare including users, advertisers and platforms is
\[
 W_r=d_rQ_r-\frac{\beta Q_r^2}{2}-n_rF_r.                \tag{5}
\]
Ad revenue $R_rQ_r$ is a transfer and is not added again to (5).
\end{theorem}
\begin{proof}
Since $b>0$ and both quantities are positive, $\CS_1>\CS_0$ is equivalent to $Q_1>Q_0$. Substitute (2) and divide positive quantities to obtain (4). The concentration claim follows from the symmetric shares.

Gross user value is $(a+h_r)Q_r-bQ_r^2/2$, and gross advertiser value is $A_rQ_r-BQ_r^2/2$. Subtracting real costs $c_rQ_r+n_rF_r$ gives (5). Alternatively, adding user surplus, advertiser surplus and net platform profits cancels user and advertising payments once each. Both calculations coincide, proving the transfer claim.
\end{proof}

Here is an exact example satisfying all entry inequalities. Set $a=3$, $A_0=A_1=2$, $c_0=c_1=1$, $b=B=1/2$, $h_0=0$, $h_1=1$, and $F_0=4/5$, $F_1=2$. Then $d_0=4,d_1=5$ and
\[
 n_0=3,\quad n_1=2,\quad Q_0=3,\quad Q_1=10/3.
\]
Indeed, baseline operating profits at three and four firms are $1$ and $16/25$, straddling $4/5$. New-regime profits at two and three firms are $25/9$ and $25/16$, straddling $2$. The user prices are $3/2$ and $7/3$, ad prices are $1/2$ and $1/3$, and
\[
 \CS_0=9/4,\quad\CS_1=25/9,\qquad
 W_0=51/10,\quad W_1=64/9.
\]
Thus concentration rises from $1/3$ to $1/2$ while both welfare measures rise and ad price falls. Consumers receive a one-unit higher privacy value for any fixed service unit; that gain is large enough to raise total use despite higher user prices. Every number is generated by the stated game, not an empirical estimate.

A fall in ad price by itself therefore does not sign consumer welfare. For example holding $A_r$ fixed, any output increase mechanically lowers $R_r$ while increasing advertiser surplus. Conversely a tracking restriction can lower $A_r$ and reduce $Q_r$, with an ambiguous net ad-price effect $\Delta R=\Delta A-B\Delta Q$.

\section{Welfare sensitivity when privacy valuations and fixed costs are unknown}
Let $h_r^w$ denote the normative privacy valuation per unit and define $e_r=h_r^w-h_r$. The price game continues to use behavioral $h_r$. For a stated evaluator, welfare becomes
\[
 W_r^w=d_rQ_r-\beta Q_r^2/2+e_rQ_r-n_rF_r.              \tag{6}
\]
This adjustment can represent better informed harm valuation. It is not identified by observed acceptance of tracking without assumptions connecting behavior to informed utility.

Suppose $b,B,Q_r,n_r$ are known and the quantity model fits. Equation (2) identifies
$d_r=\beta(n_r+1)Q_r/n_r$. The selected free-entry count identifies fixed costs only within
\[
 \pi_r(n_r+1)<F_r\leq\pi_r(n_r).                        \tag{7}
\]
Assume $e_r\in[\underline e_r,\overline e_r]$ and that all such valuations and fixed costs satisfying (7) are allowed independently across regimes. Set
$V_r=d_rQ_r-\beta Q_r^2/2$.

\begin{proposition}[Sharp interval with explicit endpoint attainment]
The sharp set for regime-$r$ welfare under these restrictions is
\[
 [L_r,U_r),\quad
 L_r=V_r+\underline e_rQ_r-n_r\pi_r(n_r),\quad
 U_r=V_r+\overline e_rQ_r-n_r\pi_r(n_r+1).                \tag{8}
\]
Consequently the sharp set for $W_1^w-W_0^w$ is
\[
 (L_1-U_0,\ U_1-L_0),                                  \tag{9}
\]
and its closure is the interval with those endpoints. Welfare is robustly nonnegative whenever $L_1-U_0\geq0$, and robustly negative whenever $U_1-L_0\leq0$.
\end{proposition}
\begin{proof}
In (6), welfare increases continuously with $e_r$ because $Q_r>0$ and decreases continuously with $F_r$ because $n_r>0$. Its minimum is attained at $(\underline e_r,\pi_r(n_r))$. Its supremum occurs as $e_r=\overline e_r$ and $F_r$ decreases to $\pi_r(n_r+1)$, which is excluded by the maximal-entry selection in (3); hence that upper endpoint is unattained. Every intermediate value is attained because the rectangular parameter domain is convex and its linear image is an interval; concretely one can vary continuously along a segment from the minimizer to a point arbitrarily close to the supremum. All these parameter vectors preserve the observed quantities, prices and entry counts under the stated model, proving sharpness. Subtracting two independently attainable half-open intervals gives precisely (9): the lower endpoint would require the unattainable $U_0$, and the upper would require the unattainable $U_1$. The sign claims follow directly.
\end{proof}

This sharp set is conditional on fixed $b,B$ and the model's equilibrium selection. If the policy mechanically links $F_1$ and $F_0$, the independent rectangles overstate the compatible set; the welfare contrast should instead be optimized over that linked set. If costs are directly observed, their corresponding interval disappears. A separate external privacy harm can be subtracted from (6), but its units, population and valuation must be specified.

\section{What remains for a policy application}
The bundle $(h,A,c,F)$ admits different channels: informed user valuation shifts, loss of advertiser targeting value, marginal compliance costs and fixed entry burdens. Observing only service quantity and firm counts identifies the combined net intercept $d$, not its components. User and ad prices provide additional equations (1), but identifying $h$ still requires the baseline demand intercept and information model. The normative correction $e$ remains unrestricted by the price equations.

A credible application therefore needs a regime definition, demand and ad-demand estimates, a baseline entry cohort, and either fixed-cost evidence or the sensitivity set (8). Independent variation in privacy information, data access and compliance burden could test separate channels. This benchmark establishes coexistence and accounting, but not concentration effects of the GDPR or of any other particular regulation, and not a universal welfare ranking of stronger privacy regimes.

\section{Completion Estimate}
50\%

This is a post hoc, heuristic estimate for the full catalogue target.
The attempt solves a two-sided quantity-entry benchmark, gives an exact privacy-concentration coexistence criterion, and derives sharp cost and welfare-valuation sensitivity intervals. Full regulatory effects still require privacy differentiation, endogenous advertising and data access, credible channel variation, and empirically identified informed consumer preferences and entry costs.

\begin{thebibliography}{9}
\bibitem{jsg} G. A. Johnson, S. K. Shriver and S. G. Goldberg (2023), \emph{Privacy and Market Concentration: Intended and Unintended Consequences of the GDPR}, Management Science 69(10), 5695--5721. Publisher abstract and metadata inspected 9 October 2026; the PDF link redirected to the abstract, so full journal text is not claimed verified. \url{https://pubsonline.informs.org/doi/abs/10.1287/mnsc.2023.4709}.
\bibitem{atw} A. Acquisti, C. Taylor and L. Wagman (2016), \emph{The Economics of Privacy}, Journal of Economic Literature 54(2), 442--492. Publisher abstract and metadata inspected 9 October 2026. \url{https://www.aeaweb.org/articles?id=10.1257/jel.54.2.442}.
\end{thebibliography}

\end{document}
